\begin{frontmatter}

\title{Cheap Signals, Costly Proof: Award-Layer Triage for Cartel Enforcement\texorpdfstring{\footnote{This version: May 2026. We thank seminar participants at INSPER for comments on earlier drafts; all remaining errors are our own.}}{}}

\author[insper]{Darcio Genicolo-Martins\corref{cor1}}
\ead{darciogm1@insper.edu.br}
\author[insper]{Paulo Furquim de Azevedo}
\affiliation[insper]{organization={INSPER},
  city={S\~ao Paulo}, country={Brazil}}
\cortext[cor1]{Corresponding author.}

\begin{abstract}
Most antitrust enforcers see who won a public contract and at
what price---not what every bidder offered. We show this
thinner data layer is enough to triage which firms warrant
costly forensic investigation. We propose an enforcement
architecture in which an award-layer screening stage runs
before a costly bid-layer forensic stage, and instantiate it
with a \emph{frequent-loser} flag built from São Paulo's BEC
contract-award records ($\valYearStart$--$\valYearEnd$). The
flag cuts the bid-microdata pool the forensic stage must work
on by $\valArchFootprintReduction$ ($\valArchSeqKone$ of
$\valArchPoolSize$ firms) while recovering
$\valArchTPThouSeqK$ of $\valArchNPos$ adjudicated cobidders.
The screen is a triage device, not an adjudication device.
Discrimination accuracy is AUC $\valAUCFLfirmTemp$ under
temporal holdout; the screen adds non-redundant signal to the
Imhof--Wallimann bid-distribution pipeline ($+\valImhofIncFL$
AUC, DeLong $p=\valImhofIncFLp$). A simple separating-equilibrium
argument motivates endogenous loser-side participation as the
ranking primitive. Wherever award records are routinely
available while bid microdata are forensic-recoverable,
screening should be sequenced before forensics.
\end{abstract}

\end{frontmatter}

\noindent\textbf{Keywords:} screening under incomplete observability, cover bidding, cartel adjacency, award-layer enforcement, separating equilibrium, participation-based statistics

\noindent\textbf{JEL No.:} D44, D73, H57, K21, L41

\bigskip
